% SPLINCE vs PCRL — variance-constrained 60-cell head-to-head.
% Generated by scripts/build_splince_vs_pcrl.py.
\begin{table}[t]
\centering
\small
\begin{tabular}{l l c c c c c c}
\toprule
Dataset & Method & Cells & $R^2_{\text{1H}}$ & $R^2_{\text{DA}}$ & $R^2$-pass & Adj-pass & Strict \\
\midrule
\multirow{2}{*}{adult} & PCRL (R5/R7) & 8 & 0.0160 & -- & 8/8 & -- & -- \\
& SPLINCE & 8 & 0.0024 & 0.0031 & 8/8 & 5/8 & 0/8 \\
\bottomrule
\end{tabular}
\caption{SPLINCE (Holstege et al., NeurIPS 2025) applied as a closed-form post-hoc eraser on the same PCRL-trained backbone, compared against the PCRL R5/R7 strict-pass baseline. SPLINCE projections are fit per (purpose, attribute, seed) cell on (features, sensitive attribute, primary task label) per Theorem 1 of the SPLINCE paper. $R^2_{\text{1H}}$ is the linear-regression $R^2$ on the centered one-hot encoding of the sensitive attribute; $R^2_{\text{DA}}$ is the dominant-axis variant of Section~\ref{sec:dominant_axis}. $R^2$-pass requires $R^2_{\text{1H}} < 0.05$; Adj-pass adds the post-hoc auditor delta $\Delta_{\text{aud}} < 0.02$; Strict adds $\bar{\sigma}_{\text{dim}} \geq 0.5$ and effective rank $\geq 2$.}
\label{tab:splince_vs_pcrl}
\end{table}
